% =====================================================================
% 10. Reproducción
% =====================================================================
\section{Reproducción del flujo}
\label{sec:reproduccion}

\subsection{Pasos}

Todo el proyecto se reproduce desde una base vacía con:

\begin{lstlisting}
git clone https://github.com/SamuOsorio/Proyecto-DBA.git && cd Proyecto-DBA
cp .env.example .env
python3 -m venv .venv && source .venv/bin/activate
pip install -r requirements.txt
# copiar el ZIP del MGN 2025 (nivel municipio) en data/raw/mgn_2025/MGN2025.zip
make all        # = make up municipios edificios analisis
\end{lstlisting}

\campo{make analisis} corre \campo{src/analisis/run.py}, que ejecuta los pasos
en orden como \campo{python -m src.analisis.<paso>}; si uno falla, se detiene
e indica cómo retomar con \campo{--desde <paso>}.
\campo{--solo-figuras} rehace gráficas, mapas, ranking y potencial a partir de
los CSV y el GeoJSON versionados, sin necesidad de MongoDB.

\begin{table}[H]
  \centering\small
  \caption{Pasos de \campo{make analisis}, en orden.}
  \label{tab:pasos}
  \begin{tabularx}{\textwidth}{@{}l L@{}}
    \toprule
    \textbf{Paso} & \textbf{Qué produce} \\
    \midrule
    \campo{asignacion} & \campo{cod\_mpio} en los edificios; registro en \campo{bitacora\_cargas} \\
    \campo{evidencia} & \campo{explain()} con IXSCAN y con COLLSCAN \\
    \campo{validacion} & Comparación con GeoPandas en 10 municipios \\
    \campo{agregacion} & \campo{resultados}: 170 documentos por fuente y umbral, verificados \\
    \campo{sensibilidad} & Tabla de sensibilidad al umbral de Google \\
    \campo{exportar} & CSV por municipio y subregión; tabla por subregión \\
    \campo{figuras\_eda} & Las figuras del EDA \\
    \campo{graficas} & Dispersión por municipio y gráfica de sensibilidad \\
    \campo{mapas} & Mapas por fuente, de diferencia y de \% de techos \\
    \campo{ranking} & Top 15, tabla y gráfica; ranking completo en CSV \\
    \campo{potencial} & Potencia instalable y energía por municipio \\
    \bottomrule
  \end{tabularx}
\end{table}

\subsection{Corrida desde base vacía}

Para este informe se borró la base \campo{upme\_solar} y se corrió
\campo{make all}. Terminó sin errores en unos 14 minutos, incluyendo el
prefiltro y la limpieza de Google desde los archivos crudos, y reprodujo
exactamente las mismas cifras: 170 municipios, 2\,694\,786 y 1\,761\,048
edificios, el mismo edificio sin municipio y 680 documentos en
\campo{resultados}. Los CSV regenerados solo difieren en la última cifra
decimal, por redondeo de punto flotante.

\fig[0.75\textwidth]{cap_make_all.png}{0.32}{Corrida de \campo{make all}
desde una base vacía: resumen de tiempos de la etapa de análisis
(\campo{run.py}).}{fig:make-all}

\subsection{Parámetros y entorno}

Los parámetros no están escritos en el código: los umbrales de confianza, la
regla de asignación, los CRS, el tamaño de los lotes, los supuestos de la
estimación de potencia y las rutas de salida se leen de \campo{config.yaml}.
La URI de la base vive en \campo{.env}, fuera del repositorio. Cada documento
de \campo{resultados} guarda el commit del código que lo produjo.

Entorno: MongoDB 8 en Docker Compose y Python 3.12 o superior con las
versiones de \campo{requirements.txt}. El flujo se corrió en macOS y en Linux,
con las mismas cifras.

\subsection{Lecciones del proceso}

\begin{itemize}
  \item Un validador \campo{\$jsonSchema} copiado tal cual del diccionario
        bloquea el flujo si exige campos que se calculan en una etapa
        posterior; hubo que declararlos opcionales o nulos.
  \item Que una geometría sea válida en el plano (GEOS) no garantiza que lo
        sea en la esfera (S2): crear el índice \campo{2dsphere} antes de
        cargar convirtió un fallo de todo el índice en el rechazo de unos
        pocos documentos reparables.
  \item Separar cada tarea en su propio archivo, comunicadas por archivos
        intermedios y colecciones, permitió que cuatro personas trabajaran en
        paralelo, a cambio de coordinar el orden de los pasos en
        \campo{run.py}.
  \item Guardar los CSV de resultados en el repositorio permitió rehacer
        mapas y figuras sin tener la base cargada, lo que aceleró la revisión
        de las gráficas.
\end{itemize}
